\begin{abcliste}
\abc
Die Kräfte (Schritt 1) und die Gleichung für Kiste 1 (Schritt 2) stimmen. \result{\text{Schritt 3 ist falsch.}} Kiste 2 ruht nicht, sie wird nach unten beschleunigt. Dann können sich die Kräfte auf sie nicht aufheben: Die Gewichtskraft muss grösser sein als die Seilkraft, und ihre Differenz ergibt $m_2\,a$. Die richtige Gleichung für Kiste 2 (nach unten positiv) lautet
\[
 m_2\,g - F_\mathrm{S} = m_2\,a
\]
Schritt 4 übernimmt den Fehler.

\abc
Aus Schritt 2, $F_\mathrm{S} = m_1\,a$, eingesetzt in die Gleichung für Kiste 2: $m_2\,g - m_1\,a = m_2\,a$, also
\begin{align}
 a &= \aKF \\
   &= \frac{\mZ\cdot\g}{\mE+\mZ} \\
   &= \aK \approx \result{\aKP{0}{2}}
\end{align}
Dasselbe folgt direkt für beide Kisten zusammen: Die Gewichtskraft von Kiste 2 beschleunigt die Masse $m_1 + m_2$; die Seilkräfte sind innere Kräfte.

\abc
\begin{align}
 F_\mathrm{S} &= m_1\,a = \FKF \\
   &= \mE\cdot\aK \\
   &= \FK \approx \result{\FKP{0}{3}}
\end{align}
Die Seilkraft ist kleiner als die Gewichtskraft von Kiste 2, $m_2\,g = \FSP{0}{3}$: Das Seil hält Kiste 2 nicht im Gleichgewicht, sonst würde sie nicht beschleunigt. Nur wenn die Kisten ruhten (oder sich mit konstanter Geschwindigkeit bewegten), wäre $F_\mathrm{S} = m_2\,g$.
\end{abcliste}
